\chapter{Representaciones unitarias y su forma integrada}\label{cap:v4-13}
Una traslación puede mirarse como un movimiento del grupo o como un operador sobre funciones. Una representación convierte la primera descripción en la segunda, preservando el producto y la geometría de un espacio de Hilbert. Estudiaremos sus subespacios invariantes y después construiremos, en ambas direcciones, el enlace con el álgebra de convolución. La medida de Haar y las identidades aproximadas del capítulo anterior son las herramientas que hacen posible esa correspondencia.

\section{Continuidad, equivalencia y representaciones regulares}\label{lectura:v4-23}
En este capítulo $G$ es localmente compacto y Hausdorff, $dx$ es Haar izquierda y los espacios de Hilbert son complejos. El producto interno es lineal en el primer argumento.
\begin{definicion}
Una \emph{representación unitaria fuertemente continua}\index{representacion@representación!unitaria} es un homomorfismo $\pi:G\to\mathcal U(H)$ tal que $x\mapsto\pi(x)u$ es continua para cada $u\in H$. En particular,
\[
\pi(e)=I,\qquad \pi(x^{-1})=\pi(x)^{-1}=\pi(x)^*.
\]
Un \emph{entrelazador}\index{entrelazador} entre $\pi_1$ en $H_1$ y $\pi_2$ en $H_2$ es un operador acotado $T:H_1\to H_2$ con $T\pi_1(x)=\pi_2(x)T$ para todo $x$. Las representaciones son equivalentes si admiten un entrelazador unitario.
\end{definicion}
La topología fuerte y la débil de operadores inducen la misma topología \emph{sobre} $\mathcal U(H)$. En efecto, si $T_\alpha\to T$ débilmente y tanto $T_\alpha$ como $T$ son unitarios,
\[
\norm{(T_\alpha-T)u}^2=2\norm u^2-2\Rea\ip{T_\alpha u}{Tu}\longrightarrow0.
\]
La implicación inversa es Cauchy--Schwarz. Es esencial que el límite también esté en $\mathcal U(H)$; el argumento no afirma que todo límite débil de unitarios sea unitario.

\begin{proposicion}[Representaciones regulares]\label{v4:rep-regulares}
En $L^2(G)$ las fórmulas
\[
[\lambda(x)f](y)=f(x^{-1}y),\qquad
[\rho(x)f](y)=\Delta(x)^{1/2}f(yx)
\]
definen representaciones unitarias fuertemente continuas. Son equivalentes mediante la involución unitaria
\[
(Jf)(y)=\Delta(y)^{-1/2}f(y^{-1}),\qquad J\lambda(x)J=\rho(x).
\]
\end{proposicion}
\begin{proof}
Las leyes de homomorfismo se verifican sustituyendo las fórmulas y usando que $\Delta$ es multiplicativa. La invariancia izquierda da $\norm{\lambda(x)f}_2=\norm f_2$; la identidad $\int h(yx)dy=\Delta(x)^{-1}\int h$ da la misma igualdad para $\rho(x)$. Sus inversas corresponden a $x^{-1}$, de modo que son sobreyectivas y unitarias.

Para $f\in\Cc(G)$, las traslaciones próximas a $e$ tienen soportes en un compacto común y convergen uniformemente a $f$. La norma $L^2$ de la diferencia queda acotada por la norma uniforme multiplicada por la raíz de la medida de ese compacto. Para $\rho$ usamos además $\Delta(x)^{1/2}\to1$. La densidad de $\Cc(G)$ y la estimación $\norm{\pi(x)f-f}_2\leq2\norm{f-h}_2+\norm{\pi(x)h-h}_2$ extienden la continuidad a todo $f\in L^2$. La ley de grupo la traslada de $e$ a cualquier punto.

La fórmula de inversión de Haar prueba
\[
\norm{Jf}_2^2=\int_G\Delta(y)^{-1}|f(y^{-1})|^2dy
=\int_G|f(z)|^2dz.
\]
Como $J^2=I$, $J$ es unitario. Finalmente,
\[
[J\lambda(x)Jf](y)
=\Delta(y)^{-1/2}\Delta(x^{-1}y^{-1})^{-1/2}f(yx)
=\Delta(x)^{1/2}f(yx).
\]
\end{proof}
\begin{observacion}[Corrección de la lectura 23]
La multiplicación por $\Delta^{1/2}$, sin inversión de la variable, no proporciona el entrelazador anunciado en la fuente. Tampoco $f(y)\mapsto f(y^{-1})$ es unitario con Haar izquierda en un grupo general. La fórmula de $J$ incorpora ambos cambios y se reduce a la inversión cuando el grupo es unimodular.
\end{observacion}

\section{Subespacios invariantes y descomposición cíclica}
Un subespacio $M$ es $\pi$-invariante si $\pi(x)M\subseteq M$ para todo $x$. Aplicar también $x^{-1}$ muestra que en realidad $\pi(x)M=M$. Su clausura es invariante por continuidad de cada operador. Si $M$ es cerrado, la restricción es una representación en $M$ y $M^\perp$ también es invariante, pues
\[
\ip{\pi(x)v}{m}=\ip{v}{\pi(x^{-1})m}=0\quad(v\in M^\perp,\ m\in M).
\]
Por tanto, la proyección $P_M$ conmuta con todos los $\pi(x)$: basta comparar ambos operadores sobre $M$ y sobre $M^\perp$. Recíprocamente, si $P_M$ conmuta, para $m=P_Mm$ se tiene $\pi(x)m=P_M\pi(x)m\in M$.

\begin{definicion}
Una representación en $H\ne\{0\}$ es \emph{irreducible}\index{representacion@representación!irreducible} si sus únicos subespacios \emph{cerrados} invariantes son $0$ y $H$. El subespacio cíclico generado por $u$ es
\[
H_u=\overline{\spanop\{\pi(x)u:x\in G\}}.
\]
La representación es \emph{cíclica}\index{representacion@representación!ciclica@cíclica} si $H_u=H$ para algún $u$. La clausura es parte de esta definición: una igualdad de clausuras no equivale a que cada vector sea una combinación lineal finita.
\end{definicion}
El espacio $H_u$ es invariante porque $\pi(y)\pi(x)u=\pi(yx)u$, primero en el span y después en su clausura. Para una familia de representaciones $(\pi_j,H_j)$, su suma directa actúa coordenada a coordenada sobre la suma directa hilbertiana. Preserva la norma y es sobreyectiva. También es fuertemente continua: dado $v=(v_j)$ y $\varepsilon>0$, elegimos un conjunto finito $F$ con $\sum_{j\notin F}\norm{v_j}^2<\varepsilon^2/16$. La contribución de esas coordenadas a $\norm{\pi(x)v-\pi(x_0)v}$ es menor que $\varepsilon/2$, uniformemente en $x$; la de $F$ tiende a cero por continuidad de sus finitas representaciones. Este argumento cubre familias no numerables, pues todo vector tiene a lo sumo un conjunto numerable de coordenadas no nulas.

\begin{teorema}\label{v4:rep-ciclicas}
Toda representación unitaria es una suma directa ortogonal de representaciones cíclicas.
\end{teorema}
\begin{proof}
Ordenemos por inclusión las familias de subespacios cíclicos no nulos, cerrados y ortogonales dos a dos. La familia vacía pertenece a este conjunto. La unión de una cadena sigue siendo una familia de ese tipo: dos de sus miembros pertenecen a una misma familia de la cadena. El lema de Zorn proporciona una familia maximal $(H_j)$. Sea $M$ la clausura de su suma algebraica, identificada isométricamente con $\bigoplus_jH_j$. Es invariante por lo anterior, y su complemento ortogonal también lo es. Si $M\ne H$, un vector $0\ne u\in M^\perp$ genera un subespacio cíclico no nulo contenido en $M^\perp$. Añadirlo contradice maximalidad. Así, $M=H$, y la identificación de la suma directa entrelaza las representaciones.
\end{proof}

\section{El lema de Schur y sus consecuencias}
\begin{teorema}[Lema de Schur]\label{v4:schur}\index{Schur, lema de}
Una representación unitaria $\pi$ es irreducible si y solo si su conmutante es $\C I$. Entre dos representaciones irreducibles, todo entrelazador no nulo es un múltiplo de un unitario. Si existe uno, fijado un entrelazador unitario $V$, todos los entrelazadores son $cV$; si no son equivalentes, el único entrelazador es cero.
\end{teorema}
\begin{proof}
Si $M$ es cerrado invariante no trivial, $P_M$ pertenece al conmutante y no es escalar. Para la dirección contraria supongamos irreducibilidad. Si $T$ conmuta con la representación, tomar adjuntas de la relación para $x^{-1}$ muestra que también $T^*$ conmuta. Entonces $A=(T+T^*)/2$ y $B=(T-T^*)/(2i)$ son autoadjuntos del conmutante. Basta demostrar que cada uno es escalar.

Si $A$ no fuera escalar, su espectro tendría mínimo $m$ y máximo $M$ distintos: para un autoadjunto, espectro reducido a $\{a\}$ implicaría $\norm{A-aI}=0$ por cálculo funcional. Elijamos $m<t<M$. Las funciones $r(s)=\max(s-t,0)$ y $q(s)=\max(t-s,0)$ dan operadores $r(A),q(A)$ no nulos con producto cero. Conmutan con $\pi(x)$, porque son límites en norma de polinomios en $A$. El espacio $\overline{\operatorname{ran}r(A)}$ es entonces cerrado, invariante y no nulo. Es propio, ya que la imagen no nula de $q(A)$ le es ortogonal: $\ip{r(A)u}{q(A)v}=\ip{u}{r(A)q(A)v}=0$. Contradice irreducibilidad. Este argumento usa únicamente el cálculo funcional continuo de autoadjuntos desarrollado en el volumen III. Luego $A,B$ son escalares y también $T=A+iB$.

Sea ahora $T\pi(x)=\pi'(x)T$. Tomando adjuntas para $x^{-1}$ obtenemos $\pi(x)T^*=T^*\pi'(x)$. Por ello $T^*T=\alpha I$ y $TT^*=\beta I$, con $\alpha,\beta\geq0$. Si $T\ne0$, ambas son positivas y la igualdad $(TT^*)T=T(T^*T)$ da $\beta=\alpha$. Por tanto $\alpha^{-1/2}T$ es unitario y entrelaza. Si $V$ es un entrelazador unitario fijo y $S$ otro entrelazador, $V^*S$ conmuta con $\pi$, luego es $cI$ y $S=cV$. Esto incluye $S=0$ sin usar su inversa.
\end{proof}
\begin{corolario}
Si $G$ es abeliano, toda representación unitaria irreducible tiene dimensión uno.
\end{corolario}
\begin{proof}
Cada $\pi(x)$ conmuta con todos los $\pi(y)$, luego es escalar por Schur. Todos los subespacios cerrados son entonces invariantes. Si $\dim H>1$, cualquier recta sería un subespacio cerrado propio no nulo, contradicción.
\end{proof}

\section{Ejercicios de entrelazado e invariancia}
\begin{ejercicio}[Lectura 23, ejercicio 1]\label{ej:v4-l23-1}
Demuestre que el conjunto $C(\pi,\pi')$ de entrelazadores es un subespacio de $\Bcal(H,H')$.
\end{ejercicio}

\begin{ejercicio}[Lectura 23, ejercicio 2]\label{ej:v4-l23-2}
Demuestre que el conmutante $C(\pi)$ es un álgebra de von Neumann.
\end{ejercicio}

\begin{ejercicio}[Lectura 23, ejercicio 3]\label{ej:v4-l23-3}
Si $M$ es cerrado e invariante, demuestre que $M^\perp$ es invariante.
\end{ejercicio}

\begin{ejercicio}[Lectura 23, ejercicio 4]\label{ej:v4-l23-4}
Para $M$ cerrado, pruebe que $M$ es invariante si y solo si $P_M\in C(\pi)$.
\end{ejercicio}


\section{Representaciones del álgebra de convolución}\label{lectura:v4-24}
La lectura 24 disponible es una reconstrucción editorial independiente, no un original recuperado de Valencia Bucio. Se integra aquí su desarrollo completo, enlazado con las identidades aproximadas del capítulo \ref{cap:v4-12}; su función es justificar los pasos que conectan las lecturas originales 23 y 25--27. La correspondencia que sigue es un resultado clásico de la teoría de representaciones \cite{reconstruccion24,folland2016}.

\begin{definicion}\label{v4-l24-d:rep}
Una representación unitaria fuertemente continua de $G$ en $H$ es una aplicación $U:G\to\mathcal U(H)$ que satisface $U(xy)=U(x)U(y)$ y para la cual $x\mapsto U(x)u$ es continua para todo $u\in H$. Una $*$-representación de $L^1(G)$ es una aplicación lineal $\rho:L^1(G)\to\Bcal(H)$ que satisface
\[
\rho(f*g)=\rho(f)\rho(g),\qquad \rho(f^*)=\rho(f)^*.
\]
Se llama no degenerada si
\[
\overline{\spanop\{\rho(f)u:f\in L^1(G),\ u\in H\}}=H.
\]
\end{definicion}

\begin{lema}\label{v4-l24-l:contractive}
Toda $*$-representación $\rho:L^1(G)\to\Bcal(H)$ es contractiva:
\[
\norm{\rho(f)}\leq\norm f_1.
\]
En particular, es continua, aun si no se incluyó continuidad en su definición.
\end{lema}
\begin{proof}
Usaremos un argumento de espectros, cuyos ingredientes recordamos. La unitarización de $A=L^1(G)$ es $\widetilde A=A\oplus\C$, con
\[
(a,z)(b,w)=(a*b+zb+wa,zw),\qquad\norm{(a,z)}=\norm a_1+|z|.
\]
Tiene identidad $(0,1)$ y es un álgebra de Banach: la multiplicación es submultiplicativa por la proposición \ref{v4-l24-p:banach}, y una sucesión de Cauchy converge en cada componente. La aplicación $\widetilde\rho(a,z)=\rho(a)+zI$ es un homomorfismo unital.

Si $\lambda(0,1)-(a,0)$ tiene inverso en $\widetilde A$, su imagen es un inverso de $\lambda I-\rho(a)$. Por tanto,
\[
\sigma_{\Bcal(H)}(\rho(a))\subseteq\sigma_{\widetilde A}((a,0)).
\]
Cuando $|\lambda|>\norm a_1$, la serie geométrica
$\lambda^{-1}\sum_{n\geq0}((a,0)/\lambda)^n$ converge en $\widetilde A$ y, multiplicando sus sumas parciales, se comprueba que es el inverso requerido. En consecuencia, el espectro de $\rho(a)$ está dentro del disco de radio $\norm a_1$.

Aplicamos esto a $a=f^**f$. El operador $\rho(a)=\rho(f)^*\rho(f)$ es positivo. Para un operador positivo acotado, el teorema espectral da
\[
\norm{\rho(f)}^2=\norm{\rho(f)^*\rho(f)}
=\sup\{|\lambda|:\lambda\in\sigma(\rho(a))\}.
\]
La inclusión anterior y la proposición \ref{v4-l24-p:banach} implican
\[
\norm{\rho(f)}^2\leq\norm{f^**f}_1\leq\norm f_1^2.
\]
Tomando raíces obtenemos la conclusión. Aquí el único resultado espectral utilizado es la igualdad entre norma y radio espectral para operadores positivos, estudiada en el volumen III; no se supone continuidad de $\rho$ para trasladar inversos.
\end{proof}

\begin{proposicion}\label{v4-l24-p:nd}
Para una $*$-representación $\rho$, son equivalentes:
\begin{enumerate}
\item $\rho$ es no degenerada.
\item Si $\rho(f)u=0$ para toda $f\in L^1(G)$, entonces $u=0$.
\item $\rho(e_U)u\to u$ para todo $u\in H$.
\end{enumerate}
\end{proposicion}
\begin{proof}
Sea $D=\spanop\{\rho(f)v\}$. Un vector $u$ es ortogonal a $D$ si y solo si
\[
0=\ip{\rho(f)v}{u}=\ip{v}{\rho(f)^*u}
\quad\text{para todos }f,v.
\]
Como un vector ortogonal a todo $v$ es cero, esto equivale a $\rho(f^*)u=0$ para toda $f$. La involución es biyectiva; por ello equivale a (ii). Además, $\overline D=H$ si y solo si $D^\perp=\{0\}$. Se obtiene (i)$\Leftrightarrow$(ii).

Supongamos (i). Si $w=\rho(f)v$, entonces
\[
\norm{\rho(e_U)w-w}
=\norm{\rho(e_U*f-f)v}\leq\norm{e_U*f-f}_1\norm v\longrightarrow0.
\]
Por linealidad vale para $w\in D$. Dado $u\in H$ y $\varepsilon>0$, tomamos $w\in D$ con $\norm{u-w}<\varepsilon/3$. Como $\norm{\rho(e_U)}\leq1$,
\[
\norm{\rho(e_U)u-u}\leq2\norm{u-w}+\norm{\rho(e_U)w-w}<\varepsilon
\]
para índices suficientemente posteriores. Así obtenemos (iii). Finalmente, (iii) implica (ii): si todas las $\rho(f)$ anulan $u$, también $\rho(e_U)u=0$, y al tomar el límite resulta $u=0$.
\end{proof}

\section{Integración de una representación}\index{representacion@representación!integrada}
\begin{proposicion}\label{v4-l24-p:integral}
Sea $U$ una representación unitaria fuertemente continua. Para cada $f\in L^1(G)$ existe un único operador $\Pi_U(f)\in\Bcal(H)$ tal que
\begin{equation}\label{v4-l24-eq:weak}
\ip{\Pi_U(f)u}{v}=\int_Gf(x)\ip{U(x)u}{v}\,dx
\quad(u,v\in H).
\end{equation}
Además, $\norm{\Pi_U(f)}\leq\norm f_1$.
\end{proposicion}
\begin{proof}
El coeficiente $x\mapsto\ip{U(x)u}{v}$ es continuo y está acotado por $\norm u\norm v$, por Cauchy--Schwarz y unitariedad. La integral de \eqref{v4-l24-eq:weak} existe absolutamente y su módulo es menor o igual que $\norm f_1\norm u\norm v$. Para $u$ fijo es una funcional antilineal continua de $v$. El teorema de representación de Riesz produce un único vector $w$ que la representa como $\ip{w}{v}$. Definimos $\Pi_U(f)u=w$.

La linealidad respecto de $u$ se verifica tomando producto interno con todo $v$: ambas expresiones para una combinación lineal tienen los mismos coeficientes y por unicidad son iguales. Asimismo,
\[
\norm{\Pi_U(f)u}=\sup_{\norm v=1}|\ip{\Pi_U(f)u}{v}|
\leq\norm f_1\norm u.
\]
Así, el operador es acotado con la norma anunciada. Si otro operador tiene los mismos coeficientes, su diferencia aplicada a $u$ es ortogonal a todo $v$ y por ello es cero. Esto prueba unicidad.
\end{proof}

Escribiremos $\Pi_U(f)=\int_Gf(x)U(x)dx$, entendiendo siempre esta integral por \eqref{v4-l24-eq:weak}. No se está suponiendo continuidad en norma de operadores. Una representación fuertemente continua puede no ser continua en esa norma. Este cuidado evita introducir una hipótesis que no pertenece al teorema.

\begin{teorema}[Representación integrada]\label{v4-l24-t:direct}
La aplicación $f\mapsto\Pi_U(f)$ de la proposición \ref{v4-l24-p:integral} es una $*$-representación no degenerada de $L^1(G)$. Además,
\begin{equation}\label{v4-l24-eq:cov}
U(x)\Pi_U(f)=\Pi_U(L_xf).
\end{equation}
\end{teorema}
\begin{proof}
La linealidad en $f$ se obtiene de \eqref{v4-l24-eq:weak} y unicidad. Para el producto, fijamos $u,v\in H$ y calculamos
\begin{align*}
\ip{\Pi_U(f)\Pi_U(g)u}{v}
&=\int_G\int_G f(x)g(y)\ip{U(xy)u}{v}\,dy\,dx\\
&=\int_G\int_G f(x)g(x^{-1}z)\ip{U(z)u}{v}\,dz\,dx\\
&=\int_G(f*g)(z)\ip{U(z)u}{v}\,dz.
\end{align*}
Para obtener la primera línea, pasamos $U(x)$ al segundo argumento como $U(x)^*$ y aplicamos \eqref{v4-l24-eq:weak} a $g$. El valor absoluto de la integral doble está acotado por $\norm f_1\norm g_1\norm u\norm v$; Fubini es aplicable. El cambio $z=xy$ usa invariancia izquierda. La igualdad de todos los coeficientes prueba multiplicatividad.

Para la involución, \eqref{v4-l24-eq:inv} da
\begin{align*}
\ip{\Pi_U(f^*)u}{v}
&=\int_G\Delta(x^{-1})\overline{f(x^{-1})}\ip{U(x)u}{v}\,dx\\
&=\int_G\overline{f(y)}\ip{U(y^{-1})u}{v}\,dy\\
&=\int_G\overline{f(y)}\ip{u}{U(y)v}\,dy
=\ip{u}{\Pi_U(f)v}
=\ip{\Pi_U(f)^*u}{v}.
\end{align*}
En la penúltima igualdad se utiliza la antilinealidad del segundo argumento. Por unicidad, $\Pi_U(f^*)=\Pi_U(f)^*$.

Por otra parte, para todo $u$,
\begin{equation}\label{v4-l24-eq:strongapprox}
\norm{\Pi_U(e_U)u-u}
\leq\int_G e_U(y)\norm{U(y)u-u}\,dy
\leq\sup_{y\in U}\norm{U(y)u-u}\longrightarrow0.
\end{equation}
La primera desigualdad se demuestra tomando coeficientes con vectores unitarios y usando la proposición \ref{v4-l24-p:integral}. El último límite es la continuidad fuerte en $e$. La proposición \ref{v4-l24-p:nd} implica no degeneración.

Finalmente,
\[
\ip{U(x)\Pi_U(f)u}{v}
=\int_G f(y)\ip{U(xy)u}{v}\,dy
=\int_G f(x^{-1}z)\ip{U(z)u}{v}\,dz.
\]
Este es el coeficiente de $\Pi_U(L_xf)$. Queda probada \eqref{v4-l24-eq:cov}.
\end{proof}

Para demostrar el recíproco necesitaremos pasar una integral a través de un operador lineal acotado. Damos el argumento para evitar un intercambio de límites injustificado. En particular, no utilizaremos un supuesto teorema de convergencia dominada para redes arbitrarias.

\section{Integrales vectoriales y reconstrucción}\index{integral!vectorial}
\begin{lema}[Integral de una función vectorial ponderada]\label{v4-l24-l:vector}
Sea $X$ un espacio de Banach, sea $F:G\to X$ continua y acotada, y sea $f\in L^1(G)$. Existe un vector, denotado $\int_Gf(x)F(x)dx$, tal que
\[
\ell\left(\int_G f(x)F(x)dx\right)=\int_G f(x)\ell(F(x))dx
\quad(\ell\in X^*).
\]
Su norma es menor o igual que $\norm f_1\sup_x\norm{F(x)}$. Si $A:X\to Y$ es lineal acotado entre espacios de Banach, entonces
\[
A\left(\int_G f(x)F(x)dx\right)=\int_G f(x)AF(x)dx.
\]
\end{lema}
\begin{proof}
Supongamos primero que $f$ se anula fuera de un compacto $K$. Por continuidad, $F(K)$ es compacto. Para cada $n$ cubrimos $K$ por una cantidad finita de abiertos relativos en cada uno de los cuales los valores de $F$ distan menos de $1/n$ de algún valor $w_j=F(x_j)$. Quitando de cada abierto los anteriores, obtenemos una partición Borel $E_1,\ldots,E_m$ de $K$. La función simple $F_n=\sum_j w_j\chi_{E_j}$ satisface $\norm{F_n-F}<1/n$ en $K$. Pongamos
\[
v_n=\sum_j w_j\int_{E_j}f(x)dx.
\]
Al usar la partición común a $F_n$ y $F_m$, la desigualdad triangular da
\[
\norm{v_n-v_m}\leq(1/n+1/m)\norm f_1.
\]
Por completitud existe $v=\lim v_n$. Para toda $\ell\in X^*$, la estimación uniforme permite pasar al límite en las integrales escalares y obtener $\ell(v)=\int f\ell(F)$. La misma estimación prueba $\norm v\leq\norm f_1\sup\norm F$. La independencia de las particiones se obtiene comparando dos aproximaciones mediante su refinamiento común. La unicidad también se sigue de que $X^*$ separa puntos, por Hahn--Banach.

Para $f$ general, elijamos $f_n\in\Cc(G)$ con $\norm{f_n-f}_1\to0$. Las integrales ya construidas satisfacen
\[
\norm{\int f_nF-\int f_mF}\leq\norm{f_n-f_m}_1\sup\norm F.
\]
Su límite existe y no depende de la sucesión elegida. Las integrales escalares convergen por la misma cota, de modo que el límite tiene la propiedad anunciada y conserva la estimación. Por último, para cada $\eta\in Y^*$, la funcional $\eta A$ pertenece a $X^*$. Aplicar la identidad ya probada a $\eta A$ muestra que los dos vectores de la última fórmula tienen los mismos valores bajo toda $\eta$. Hahn--Banach implica que son iguales.
\end{proof}

\begin{observacion}\label{v4-l24-o:convvector}
El lema anterior se aplica a $F(x)=L_xg$ como función con valores en $L^1(G)$: el lema \ref{v4-l24-l:translation} da continuidad y $\norm{L_xg}_1=\norm g_1$. En este sentido,
\begin{equation}\label{v4-l24-eq:convvector}
f*g=\int_Gf(x)L_xg\,dx\quad\text{en }L^1(G).
\end{equation}
Para comprobar la identidad sin identificar indebidamente valores puntuales de clases $L^1$, sea $q\in L^\infty(G)$. La funcional $h\mapsto\int h(y)q(y)dy$ es acotada en $L^1$. El lema \ref{v4-l24-l:vector} y Fubini muestran que su valor sobre el lado derecho de \eqref{v4-l24-eq:convvector} es
\[
\int_G\int_G f(x)g(x^{-1}y)q(y)\,dy\,dx
=\int_G(f*g)(y)q(y)\,dy.
\]
La integral absoluta es menor o igual que $\norm f_1\norm g_1\norm q_\infty$. Estas funcionales separan puntos de $L^1$: si $h\ne0$, elegir $q=\overline h/|h|$ donde $h\ne0$ y $q=0$ en otro caso da $\int hq=\norm h_1>0$. Por tanto, los dos lados de \eqref{v4-l24-eq:convvector} son iguales.
\end{observacion}

\begin{teorema}[Reconstrucción de la representación del grupo]\label{v4-l24-t:inverse}
Sea $\rho:L^1(G)\to\Bcal(H)$ una $*$-representación no degenerada. Existe una única representación unitaria fuertemente continua $U$ de $G$ tal que $\rho=\Pi_U$. Para todo $x\in G$ se cumple
\begin{equation}\label{v4-l24-eq:reconstruct}
U(x)\rho(f)u=\rho(L_xf)u,
\qquad
U(x)=\operatorname*{s-lim}_{U_0}\rho(L_xe_{U_0}).
\end{equation}
El subíndice $U_0$ designa aquí una vecindad, para distinguirlo de la representación $U$.
\end{teorema}
\begin{proof}
\textit{Construcción sobre un subespacio denso.} Sea $D=\spanop\{\rho(f)u\}$, denso por hipótesis. Para $x\in G$ fijo, definimos provisionalmente
\[
V_x\left(\sum_{j=1}^n\rho(f_j)u_j\right)
=\sum_{j=1}^n\rho(L_xf_j)u_j.
\]
Debemos demostrar que el valor no depende de la expresión del vector. La identidad
\[
(L_xe_{U_0})*f=L_x(e_{U_0}*f)
\]
se verifica directamente mediante el cambio $y=xz$ en la convolución. Por los lemas \ref{v4-l24-l:approx} y \ref{v4-l24-l:contractive},
\[
\rho(L_xe_{U_0})\sum_j\rho(f_j)u_j
\longrightarrow\sum_j\rho(L_xf_j)u_j.
\]
Si la suma inicial es cero, cada término de la red a la izquierda es cero; su límite también lo es. Dos expresiones de un mismo vector tienen diferencia cero, de modo que sus imágenes coinciden. Esto prueba que $V_x$ está bien definida y es lineal en $D$.

\textit{Extensión acotada.} Como $\norm{\rho(L_xe_{U_0})}\leq\norm{L_xe_{U_0}}_1=1$, tomando el límite anterior obtenemos $\norm{V_xw}\leq\norm w$ para $w\in D$. Para $u\in H$, elegimos $w_n\in D$ con $w_n\to u$. La desigualdad hace que $(V_xw_n)$ sea de Cauchy. Su límite no depende de la sucesión, pues $\norm{V_xw_n-V_xz_n}\leq\norm{w_n-z_n}$. Definimos $U(x)u=\lim V_xw_n$. Esta extensión es lineal, tiene norma a lo sumo uno y es única por densidad.

\textit{Ley de grupo y unitariedad.} La fórmula sobre generadores da
\[
U(x)U(y)\rho(f)u=\rho(L_xL_yf)u
=\rho(L_{xy}f)u=U(xy)\rho(f)u.
\]
Por linealidad vale en $D$, y por continuidad en $H$. Del mismo modo $U(e)=I$. Así, $U(x^{-1})$ es el inverso de $U(x)$. Ambos son contracciones; por ello
\[
\norm u=\norm{U(x^{-1})U(x)u}\leq\norm{U(x)u}\leq\norm u.
\]
Entonces $U(x)$ es una isometría sobreyectiva. Una isometría lineal preserva el producto interno por la identidad de polarización; su inversa coincide con su adjunta. Por tanto, $U(x)$ es unitaria.

\textit{Continuidad fuerte.} Para $w=\sum_j\rho(f_j)u_j\in D$,
\[
\norm{U(x)w-U(x_0)w}
\leq\sum_j\norm{L_xf_j-L_{x_0}f_j}_1\norm{u_j}\longrightarrow0
\]
cuando $x\to x_0$, por el lema \ref{v4-l24-l:translation}. Para $u\in H$, tomamos $w\in D$ próximo a $u$ y usamos unitariedad:
\[
\norm{U(x)u-U(x_0)u}
\leq2\norm{u-w}+\norm{U(x)w-U(x_0)w}.
\]
Primero hacemos pequeño el primer término y después elegimos $x$ suficientemente cerca de $x_0$. Así obtenemos continuidad para todo $u$.

\textit{Recuperación de $\rho$.} Por \eqref{v4-l24-eq:convvector}, para $f,g\in L^1(G)$ y $u\in H$,
\begin{align*}
\Pi_U(f)\rho(g)u
&=\int_G f(x)U(x)\rho(g)u\,dx\\
&=\int_G f(x)\rho(L_xg)u\,dx
=\rho\left(\int_G f(x)L_xg\,dx\right)u\\
&=\rho(f*g)u=\rho(f)\rho(g)u.
\end{align*}
El paso de $\rho$ a través de la integral es el lema \ref{v4-l24-l:vector}, aplicado al operador acotado $h\mapsto\rho(h)u$; su acotación procede del lema \ref{v4-l24-l:contractive}. La primera integral coincide con la definición débil de $\Pi_U$ al tomar coeficientes. Los operadores $\Pi_U(f)$ y $\rho(f)$ coinciden sobre $D$ y, al ser acotados, sobre $H$.

\textit{Límite fuerte y unicidad.} En $D$ ya probamos la convergencia de $\rho(L_xe_{U_0})$ a $U(x)$. La cota uniforme uno extiende la convergencia a $H$ mediante la misma estimación de aproximación usada para continuidad. Si $W$ es otra representación con $\Pi_W=\rho$, \eqref{v4-l24-eq:cov} y \eqref{v4-l24-eq:strongapprox} dan
\[
W(x)u=\lim_{U_0}W(x)\rho(e_{U_0})u
=\lim_{U_0}\rho(L_xe_{U_0})u=U(x)u.
\]
Esto prueba unicidad y concluye el teorema.
\end{proof}

\section{Entrelazado, invariancia y ciclicidad}
\begin{teorema}[Entrelazado e invariancia]\label{v4-l24-t:intertwine}
Sean $U_j:G\to\mathcal U(H_j)$ representaciones unitarias fuertemente continuas, y sean $\Pi_j$ sus representaciones integradas. Para $T\in\Bcal(H_1,H_2)$ son equivalentes
\[
TU_1(x)=U_2(x)T\quad(x\in G),
\qquad T\Pi_1(f)=\Pi_2(f)T\quad(f\in L^1(G)).
\]
Además, un subespacio cerrado $M\subseteq H$ es invariante bajo $U(G)$ si y solo si lo es bajo $\Pi_U(L^1(G))$.
\end{teorema}
\begin{proof}
Supongamos la identidad de entrelazado para $G$. Tomando coeficientes y usando $T^*$, la integral de $f(x)TU_1(x)u$ tiene los mismos coeficientes que $T\Pi_1(f)u$. Sustituir $TU_1(x)$ por $U_2(x)T$ da $\Pi_2(f)Tu$. Por unicidad de coeficientes se obtiene la identidad para $L^1(G)$.

Recíprocamente, fijemos $x\in G$. Aplicamos la identidad para funciones a $L_xe_{U_0}$ y tomamos límites fuertes del teorema \ref{v4-l24-t:inverse}. Para cada $u$,
\[
TU_1(x)u=\lim T\Pi_1(L_xe_{U_0})u
=\lim\Pi_2(L_xe_{U_0})Tu=U_2(x)Tu.
\]
La continuidad de $T$ permite el primer paso.

Si $M$ es $U$-invariante, $u\in M$ y $v\in M^\perp$, entonces $\ip{U(x)u}{v}=0$ para todo $x$; integrar da $\ip{\Pi_U(f)u}{v}=0$. Como $M$ es cerrado, $(M^\perp)^\perp=M$, de modo que $\Pi_U(f)u\in M$.

Si $M$ es invariante bajo todas las $\Pi_U(f)$, también contiene $\Pi_U(L_xe_{U_0})u$ cuando $u\in M$. Estos vectores convergen a $U(x)u$. La cerradura de $M$ implica que el límite pertenece a $M$. Esto prueba la otra dirección.
\end{proof}

\begin{corolario}\label{v4-l24-c:invariants}
La correspondencia $U\leftrightarrow\Pi_U$ conserva equivalencia unitaria, irreducibilidad y vectores cíclicos. Además,
\[
\{T:TU(x)=U(x)T\ \forall x\}
=\{T:T\Pi_U(f)=\Pi_U(f)T\ \forall f\}.
\]
\end{corolario}
\begin{proof}
La igualdad es el teorema \ref{v4-l24-t:intertwine} con $H_1=H_2$. Si el entrelazador es unitario, ese mismo teorema prueba equivalencia en ambas direcciones. La irreducibilidad significa que los únicos subespacios \emph{cerrados} invariantes son $\{0\}$ y $H$; la equivalencia de invariancia demuestra la afirmación.

Para un vector $u$ pongamos $M_G=\overline{\spanop\{U(x)u:x\in G\}}$ y $M_A=\overline{\spanop\{\Pi_U(f)u:f\in L^1(G)\}}$. Si $v\perp M_G$, todos los coeficientes $\ip{U(x)u}{v}$ son cero y su integral también; entonces $v\perp M_A$. Esto implica $M_A\subseteq M_G$. Por otro lado, $U(x)u=\lim\Pi_U(L_xe_{U_0})u\in M_A$, de donde $M_G\subseteq M_A$. Por tanto, $M_G=M_A$, y uno es $H$ si y solo si el otro lo es.
\end{proof}

\section{Representación regular, caracteres y completaciones}
\begin{ejemplo}[Representación regular]\label{v4-l24-e:regular}
La representación regular izquierda en $L^2(G)$ viene dada por
\[
[\lambda(x)u](y)=u(x^{-1}y).
\]
Es unitaria por invariancia izquierda; su continuidad fuerte se prueba como en el lema \ref{v4-l24-l:translation}, usando densidad de $\Cc(G)$ en $L^2(G)$ y la norma $L^2$. Su representación integrada satisface
\[
\Pi_\lambda(f)u=f*u,\qquad
\norm{f*u}_2\leq\norm f_1\norm u_2.
\]
Para $f,u\in\Cc(G)$, la fórmula se obtiene sustituyendo la definición de $\lambda$ en los coeficientes e intercambiando integrales absolutamente convergentes sobre compactos. Para $f\in L^1$ y $u\in\Cc$, aproximamos $f$ en $L^1$ y usamos la cota de la proposición \ref{v4-l24-p:integral}. Después extendemos por densidad en $u$ a $L^2$; esta extensión define el operador de convolución en $L^2$. Es la construcción que utilizaremos en el capítulo \ref{cap:v4-16}.
\end{ejemplo}

\begin{ejemplo}[Caracteres y Fourier]\label{v4-l24-e:characters}
Sea $\chi:G\to\T$ un homomorfismo continuo. La fórmula $U_\chi(x)z=\chi(x)z$ define una representación unitaria en $\C$. Su forma integrada es la funcional
\[
\Phi_\chi(f)=\int_G f(x)\chi(x)\,dx,
\]
que satisface $\Phi_\chi(f*g)=\Phi_\chi(f)\Phi_\chi(g)$ y $\Phi_\chi(f^*)=\overline{\Phi_\chi(f)}$. Esto justifica el paso usado en la observación 4 de la lectura 26.

Si se adopta $\widehat f(\chi)=\int f(x)\overline{\chi(x)}dx$, entonces $\widehat f(\chi)=\Phi_{\overline\chi}(f)$. La aparición o ausencia del conjugado corresponde a elegir $\chi$ o su inverso; debe mantenerse una sola convención durante un cálculo. El teorema de correspondencia da el recíproco para funcionales no nulas que preservan la involución. Para funcionales multiplicativas arbitrarias hay que demostrar además esa compatibilidad; no se debe atribuir ese paso al teorema \ref{v4-l24-t:inverse} sin comprobar su hipótesis.
\end{ejemplo}

\begin{observacion}[Alcance de la reconstrucción]
En un grupo discreto, $\delta_x$ es una función de $L^1(G)$ y $U(x)=\Pi_U(\delta_x)$ cuando se usa medida de conteo. En un grupo no discreto, la masa de Dirac $\delta_x$ es una medida, no un elemento de $L^1(G)$. El límite de \eqref{v4-l24-eq:reconstruct} reemplaza esa evaluación inexistente. Se trata de convergencia fuerte, no necesariamente de convergencia en norma de operadores.
\end{observacion}

La lectura 27 parte de la representación regular para construir un álgebra $C^*$. Conviene registrar la distinción necesaria para grupos no abelianos.

\begin{definicion}\label{v4-l24-d:cstar}
El álgebra reducida de grupo es
\[
C_r^*(G)=\overline{\Pi_\lambda(L^1(G))}^{\,\norm{\cdot}_{\Bcal(L^2(G))}}.
\]
La norma universal de grupo sobre $L^1(G)$ es
\[
\norm f_{\mathrm{univ}}=\sup_U\norm{\Pi_U(f)},
\]
donde basta tomar un representante de cada clase de representaciones cíclicas. En efecto, la descomposición del teorema \ref{v4:rep-ciclicas} convierte toda representación en una suma directa, cuya norma es el supremo de las normas de sus sumandos. Además, si $D$ es un subconjunto denso de $L^1(G)$ de cardinal $\kappa$, una representación cíclica tiene un conjunto total de cardinal a lo sumo $\max(\kappa,\aleph_0)$, formado por combinaciones racionales complejas de $\Pi_U(D)u$. Sus espacios de Hilbert pueden por ello elegirse entre los espacios $\ell^2(I)$ con cardinal acotado por ese número. Sobre dichos espacios las representaciones forman un conjunto, de modo que el supremo no recorre una clase propia. La completación en esa norma se denota $C^*(G)$.
\end{definicion}

\begin{observacion}
La cota $\norm{\Pi_U(f)}\leq\norm f_1$ hace finito el supremo. La representación regular es fiel sobre $L^1(G)$: se demuestra en el ejercicio \ref{ej:v4-l24-5}. Por ello el supremo es una norma y no solo una seminorma. Las identidades de representación dan
\[
\norm{f^**f}_{\mathrm{univ}}
=\sup_U\norm{\Pi_U(f)^*\Pi_U(f)}
=\norm f_{\mathrm{univ}}^2.
\]
Así, su completación es un álgebra $C^*$. Una representación integrada se extiende de manera única a $C^*(G)$: si $f_n$ es de Cauchy en norma universal, la desigualdad $\norm{\Pi_U(f_n)-\Pi_U(f_m)}\leq\norm{f_n-f_m}_{\mathrm{univ}}$ proporciona el límite, y garantiza independencia y continuidad.

Para grupos localmente compactos abelianos, el teorema de Gelfand y el análisis de Fourier identifican ambas completaciones con $C_0(\widehat G)$. Este resultado pertenece al desarrollo posterior de las lecturas 25--27 y no se utiliza en las demostraciones anteriores. En general, la clausura de la representación regular debe llamarse $C_r^*(G)$; la correspondencia con \emph{todas} las representaciones unitarias se formula para $C^*(G)$.
\end{observacion}

\section{Ejercicios de la correspondencia}

Los ocho ejercicios siguientes verifican consecuencias y casos concretos del teorema. Sus soluciones forman parte del apunte y no requieren completar argumentos omitidos.

\begin{ejercicio}[Lectura 24 (reconstrucción), ejercicio 1]\label{ej:v4-l24-1}
\textbf{La representación trivial y una representación no fiel.}
Demuestre que $\varepsilon(f)=\int_G f(x)dx$ es una $*$-representación no degenerada en $\C$. Para $G=\R$, encuentre una función no nula en su núcleo.

\end{ejercicio}


\begin{ejercicio}[Lectura 24 (reconstrucción), ejercicio 2]\label{ej:v4-l24-2}
\textbf{Reconstrucción en $\mathbb Z$.}
Sea $T$ un operador unitario en $H$. Para $a=(a_n)_{n\in\mathbb Z}\in\ell^1(\mathbb Z)$ defina $\rho_T(a)=\sum_{n\in\mathbb Z}a_nT^n$. Pruebe que es una $*$-representación no degenerada y recupere su representación de grupo.

\end{ejercicio}


\begin{ejercicio}[Lectura 24 (reconstrucción), ejercicio 3]\label{ej:v4-l24-3}
\textbf{Una identidad aproximada explícita en $\R$.}
Con medida de Lebesgue, sea $e_n(x)=\tfrac n2\chi_{[-1/n,1/n]}(x)$. Demuestre que $e_n*f\to f$ en $L^1(\R)$ y que $\Pi_U(e_n)u\to u$ para toda representación unitaria fuertemente continua de $\R$.

\end{ejercicio}


\begin{ejercicio}[Lectura 24 (reconstrucción), ejercicio 4]\label{ej:v4-l24-4}
\textbf{Evaluación de Fourier como representación.}
Para $\xi\in\R$ ponga $\chi_\xi(x)=e^{-i\xi x}$ y $\Phi_\xi(f)=\int_\R f(x)e^{-i\xi x}dx$. Verifique multiplicatividad, compatibilidad con involución y $\norm{\Phi_\xi}=1$.

\end{ejercicio}


\begin{ejercicio}[Lectura 24 (reconstrucción), ejercicio 5]\label{ej:v4-l24-5}
\textbf{Fidelidad de la representación regular en $L^1(G)$.}
Si $\Pi_\lambda(f)=0$ como operador sobre $L^2(G)$, demuestre que $f=0$ en $L^1(G)$. Explique por qué esto no identifica automáticamente $C^*(G)$ con $C_r^*(G)$.

\end{ejercicio}


\begin{ejercicio}[Lectura 24 (reconstrucción), ejercicio 6]\label{ej:v4-l24-6}
\textbf{Por qué es indispensable la no degeneración.}
Sea $U_0$ una representación en $H_0$, sea $K\ne\{0\}$ y defina $\rho(f)=\Pi_{U_0}(f)\oplus0$ en $H_0\oplus K$. Pruebe que $\rho$ es una $*$-representación que no puede ser la representación integrada de una representación unitaria en todo $H_0\oplus K$.

\end{ejercicio}


\begin{ejercicio}[Lectura 24 (reconstrucción), ejercicio 7]\label{ej:v4-l24-7}
\textbf{Conmutante e irreducibilidad.}
Pruebe que $U$ es irreducible si y solo si todo operador que conmuta con todas las $\Pi_U(f)$ es un múltiplo escalar de $I$. Precise qué significa aquí ``invariante''.

\end{ejercicio}


\begin{ejercicio}[Lectura 24 (reconstrucción), ejercicio 8]\label{ej:v4-l24-8}
\textbf{Convergencia fuerte que no es convergencia en norma.}
Para la representación regular de $\R$, pruebe que los operadores del ejercicio \ref{ej:v4-l24-3} convergen fuertemente a $I$, pero no en norma de operadores. Puede utilizar el teorema de Plancherel y la convención $\widehat f(\xi)=\int_\R f(x)e^{-i\xi x}dx$.

\end{ejercicio}




\paragraph{Cierre.} La representación del grupo y su forma integrada contienen la misma información cuando se exige no degeneración. La continuidad fuerte permite recuperar cada operador del grupo mediante funciones integrables concentradas cerca de un punto; no hace falta convertir una masa de Dirac en una función. En los capítulos siguientes, los caracteres serán el caso unidimensional de esta correspondencia y la representación regular proporcionará el paso de Fourier a las álgebras $C^*$.
